\documentclass[aspectratio=169,11pt]{beamer}

\usetheme{Madrid}
\usecolortheme{default}
\usefonttheme{professionalfonts}
\setbeamertemplate{navigation symbols}{}
\setbeamertemplate{footline}[frame number]

\usepackage{amsmath, amssymb, booktabs, array}

\title{Synthesis, Welfare, and Student Research Designs}
\subtitle{Behavioral Labor -- Week 10}
\author{Teaching deck draft}
\date{}

\begin{document}

\begin{frame}
  \titlepage
\end{frame}

\begin{frame}{Central question}
\begin{itemize}
    \item How do we translate behavioral labor ideas into welfare-relevant, empirically credible, frontier research projects?
    \item Week 10 is a capstone studio, not a review lecture.
    \item The goal is to move from a labor-market puzzle to a wedge, setting, design, method, welfare interpretation, and contribution.
    \item A good project integrates workers, workplaces, firms, markets, and policy design.
\end{itemize}
\end{frame}

\begin{frame}{Why a final synthesis week matters}
\begin{itemize}
    \item Students can know many papers and still lack a research design.
    \item Behavioral labor claims are easy to overstate if treatment effects, mechanisms, welfare, and equilibrium are blurred.
    \item The field is strongest when behavioral objects are tied to labor-market margins and institutions.
    \item The final deliverable is a draftable research memo.
\end{itemize}
\end{frame}

\begin{frame}{Course architecture in one map}
\small
\begin{tabular}{p{0.22\linewidth} p{0.28\linewidth} p{0.34\linewidth}}
\toprule
Course block & Behavioral object & Research-design task \\
\midrule
Weeks 1--4 & preferences, beliefs, attention, learning & define the wedge and labor margin \\
Weeks 5--6 & incentives, reciprocity, identity, norms & embed the wedge inside workplaces and social settings \\
Week 7 & experiments, measurement, identification & separate treatment effects from behavioral objects \\
Weeks 8--9 & firm response, markets, policy & add equilibrium, implementation, and welfare \\
Week 10 & synthesis and project design & build a credible frontier question \\
\bottomrule
\end{tabular}
\end{frame}

\begin{frame}{What counts as a behavioral-labor contribution?}
\[
Y = g(B, I, M, P, X; \theta)
\]

\small
\begin{tabular}{p{0.12\linewidth} p{0.70\linewidth}}
\toprule
Object & Interpretation \\
\midrule
$B$ & behavioral wedge: beliefs, attention, present bias, fairness, identity, norms, learning \\
$I$ & institution or environment: tax schedule, contract, benefit rule, workplace, platform \\
$M$ & firm or market response: menus, monitoring, sorting, exploitation, accommodation \\
$P$ & policy or program design: defaults, simplification, information, assistance, targeting \\
$Y$ & labor outcome: work, effort, search, take-up, saving, training, sorting, earnings \\
\bottomrule
\end{tabular}
\end{frame}

\begin{frame}{A contribution needs all four pieces}
\begin{itemize}
    \item A real labor margin: search, effort, hours, saving, training, take-up, sorting, retention, promotion.
    \item A specific behavioral wedge, not a generic ``friction.''
    \item An institution where the wedge changes the labor prediction.
    \item A design that identifies the object being claimed.
    \item A bounded interpretation: treatment effect, behavioral parameter, welfare wedge, or equilibrium response.
\end{itemize}
\end{frame}

\begin{frame}{Welfare and interpretation}
\[
W = \int \Psi_i\big(a_i, a_i^\star, M, P\big)\, dF(i)
\]

\begin{itemize}
    \item Observed choice may not fully reveal welfare when choices are distorted.
    \item The hard object is $a_i^\star$: full-information action, commitment-consistent action, long-run learned action, or model benchmark.
    \item Defaults, take-up, search beliefs, and attention interventions all require explicit normative criteria.
    \item A large behavioral response is not automatically a welfare gain.
\end{itemize}
\end{frame}

\begin{frame}{Topic $\rightarrow$ wedge $\rightarrow$ setting}
\small
\begin{tabular}{p{0.25\linewidth} p{0.26\linewidth} p{0.31\linewidth}}
\toprule
Topic family & Candidate wedge & Labor setting \\
\midrule
Labor supply, saving, training & present bias, commitment, reference points & daily work, payroll saving, training enrollment \\
Job search and beliefs & biased expectations, slow learning & unemployment spells, applications, wage expectations \\
Attention and salience & perceived schedules, complexity & taxes, transfers, claiming, nonlinear contracts \\
Workplace behavior & reciprocity, fairness, monitoring response & effort, output, attendance, subjective evaluation \\
Identity, norms, culture & social image, norm compliance & sorting, retention, team behavior, promotion \\
Policy and firms & defaults, burden, targeting, market response & benefit design, menus, platforms, employer policy \\
\bottomrule
\end{tabular}
\end{frame}

\begin{frame}{Estimand versus interpretation}
\[
\tau = \mathbb{E}[Y_i(1)-Y_i(0) \mid S_i=1]
\]

\begin{itemize}
    \item Sometimes $\tau$ is exactly the object: did the intervention change the labor outcome?
    \item Sometimes the target is a behavioral parameter: belief bias, present bias, attention, reciprocity.
    \item Sometimes the target is welfare: corrected choice, targeting, or benchmark utility.
    \item Sometimes the target is equilibrium: firm response, spillovers, pass-through, or market adaptation.
\end{itemize}
\end{frame}

\begin{frame}{Mapping setting to design and method}
\scriptsize
\begin{tabular}{p{0.30\linewidth} p{0.31\linewidth} p{0.23\linewidth}}
\toprule
Setting & Common methods & Diagnostic question \\
\midrule
Commitment or labor-supply menus & field experiments, treatment contrasts, structural recovery & does the design isolate self-control or schedule perception? \\
Repeated job-search beliefs & panel FE, forecast errors, hazards & do beliefs predict or respond to search outcomes? \\
Nonlinear schedules and knowledge & bunching, local elasticities, learning panels & statutory or perceived incentive? \\
Workplace reciprocity and monitoring & field experiments, heterogeneity, interactions & incentive response or social-preference channel? \\
Take-up and implementation & encouragement, hazard timing, admin data & claiming value or burden reduction? \\
Welfare and targeting & sufficient statistics, calibration, structure & what benchmark action is respected? \\
\bottomrule
\end{tabular}
\end{frame}

\begin{frame}{When reduced form is enough}
\begin{itemize}
    \item The project asks whether a clearly defined intervention changes a labor outcome.
    \item The behavioral channel is sharply embedded in treatment design.
    \item The contribution is local but meaningful: take-up, effort, search, training, saving, or claiming.
    \item Welfare and equilibrium claims are stated as limits or supported by auxiliary evidence.
\end{itemize}
\end{frame}

\begin{frame}{When structure is needed}
\[
\hat \psi \in \arg\min_{\psi \in \Psi}
\left[m^{data}-m^{model}(\psi)\right]'W\left[m^{data}-m^{model}(\psi)\right]
\]

\begin{itemize}
    \item The project needs unobserved parameters: present bias, learning, search costs, social preferences.
    \item The project makes counterfactual policy or targeting claims outside the treatment support.
    \item Welfare depends on a distorted-choice benchmark that is not directly observed.
    \item Equilibrium response or firm adaptation is central to interpretation.
\end{itemize}
\end{frame}

\begin{frame}{Frontier project families}
\begin{itemize}
    \item Dynamic learning of policy schedules, eligibility, claiming, and training rules.
    \item Belief formation in rapidly changing labor markets.
    \item Algorithmic interfaces in search, scheduling, productivity, and benefits.
    \item Behavioral contracting, monitoring, feedback, and subjective evaluation inside firms.
    \item Behavioral heterogeneity, targeting, and distributional welfare.
    \item Firm-side exploitation, accommodation, insurance, and sorting around worker biases.
\end{itemize}
\end{frame}

\begin{frame}{Why the frontier is still open}
\begin{itemize}
    \item Behavioral objects are hard to measure at scale.
    \item Worker response may mix beliefs, attention, constraints, trust, and stigma.
    \item Short-run treatment effects can decay, persist, diffuse, or reverse through learning.
    \item Firms and platforms adapt to predictable worker behavior.
    \item Welfare requires a normative benchmark, not just a coefficient.
\end{itemize}
\end{frame}

\begin{frame}{Common failure modes}
\begin{itemize}
    \item Behavior is measured, but the behavioral wedge is not.
    \item The wedge is identified, but the labor-market margin is weak.
    \item A reduced-form effect is interpreted as a structural parameter.
    \item Learning dynamics are ignored.
    \item Firm or market response is offstage but central.
    \item A welfare claim is made without a benchmark.
    \item The project is interesting but not distinctly behavioral or not distinctly labor.
\end{itemize}
\end{frame}

\begin{frame}{Research memo template}
\[
Contribution = Question + Mechanism + Data + Design + Interpretation
\]

\small
\begin{tabular}{p{0.25\linewidth} p{0.56\linewidth}}
\toprule
Memo component & What it must say \\
\midrule
Question & labor-market puzzle or fact \\
Mechanism & behavioral wedge and standard alternative \\
Data & observed labor margin and behavioral object \\
Design & identifying variation, method, and estimand \\
Interpretation & welfare, policy, firm response, or equilibrium scope \\
\bottomrule
\end{tabular}
\end{frame}

\begin{frame}{Week 10 lab as proposal studio}
\begin{itemize}
    \item Reproduce: reverse-engineer one or two anchor papers' question and design logic.
    \item Diagnose: name the behavioral object, labor margin, method, welfare claim, and limit.
    \item Transfer: develop a mini-project memo in a new setting.
    \item Anchor menu: job-search beliefs, workplace social preferences, take-up frictions, retirement defaults.
\end{itemize}
\end{frame}

\begin{frame}{Bridge outward}
\begin{itemize}
    \item A dissertation-quality project usually starts as a disciplined memo.
    \item The frontier rewards sharper objects, not broader lists of anomalies.
    \item The strongest projects connect worker behavior, firm response, policy design, and welfare.
    \item Behavioral Labor should leave students with open questions and research traction.
\end{itemize}
\end{frame}

\begin{frame}{Takeaways}
\begin{itemize}
    \item Start from a labor margin and name the behavioral wedge.
    \item Embed the wedge in an institution, firm, market, or policy environment.
    \item Match the empirical method to the object being identified.
    \item Treat welfare and equilibrium as interpretation choices that need evidence.
    \item Turn synthesis into a research-design memo.
\end{itemize}
\end{frame}

\end{document}
